% Part: first-order-logic
% Chapter: model-theory
% Section: substructures

\documentclass[../../../include/open-logic-section]{subfiles}

\begin{document}

\olfileid{mod}{bas}{sub}
\olsection{উপ\printtoken{p}{structure}}

কোনো !!a{structure}~$\Struct{M}$-এর !!{domain} অন্য
গঠন~$\Struct{M'}$-এর সংজ্ঞাক্ষেত্রের উপসেট হতে পারে। কিন্তু
$\Struct{M}$-কে স্পষ্টতই $\Struct{M'}$-এর একটি “অংশ” বলা উচিত কেবল তখনই,
যখন শুধু $\Domain{M} \subseteq \Domain{M'}$ নয়, অন্তত তাদের অভিন্ন
অংশ~$\Domain{M}$-এ ভাষার সংকেতগুলি কীভাবে ব্যাখ্যা করা হয়, সে বিষয়েও
$\Struct{M}$ ও~$\Struct{M'}$ “একমত”।

\begin{defn}
\ollabel{defn:substructure}
একই ভাষা~$\Lang L$-এর !!{structure}s~$\Struct M$ ও~$\Struct M'$ দেওয়া
থাকলে, $\Struct M$-কে $\Struct M'$-এর \emph{উপ!!{structure}} এবং
$\Struct M'$-কে $\Struct M$-এর \emph{সম্প্রসারণ} বলি, আর লিখি
$\Struct M \substruct \Struct M'$, যদি এবং কেবল যদি
\begin{enumerate}
\item $\Domain{M} \subseteq \Domain{M'}$;
\item প্রত্যেক ধ্রুবক $c \in \Lang L$-এর জন্য $\Assign{c}{M} =
    \Assign{c}{M'}$;
\item প্রত্যেক $n$-স্থানীয় !!{function} $f \in \Lang L$-এর জন্য
  সব $a_1$, \dots, $a_n \in \Domain{M}$-এর ক্ষেত্রে
  $\Assign{f}{M}(a_1, \dots, a_n) = \Assign{f}{M'}(a_1, \dots, a_n)$;
\item প্রত্যেক $n$-স্থানীয় !!{predicate} $R \in \Lang L$-এর জন্য,
  সব $a_1$, \dots, $a_n \in \Domain{M}$-এর ক্ষেত্রে $\langle
  a_1, \dots, a_n\rangle \in \Assign{R}{M}$ যদি এবং কেবল যদি
  $\langle a_1, \dots, a_n\rangle \in \Assign{R}{M'}$।
\end{enumerate}
\end{defn}

\begin{rem}
\ollabel{rem:substructure}
ভাষাটিতে কোনো ধ্রুবক বা !!{function}s না থাকলে, যেকোনো অশূন্য $N
\subseteq \Domain{M}$ সেট~$\Struct M$-এর একটি
উপ!!{structure}~$\Struct{N}$ নির্ধারণ করে,
যার !!{domain}~$\Domain{N} = N$; এর জন্য $\Assign{R}{N} =
\Assign{R}{M} \cap N^n$ ধরা হয়।
% BN-SRC-112 TARGET-CORRECTION-BEGIN
\par\smallskip\noindent\textnormal{\small\textbf{উৎস-সংশোধন (BN-SRC-112):} হিমায়িত ইংরেজি মন্তব্যটিতে যেকোনো $N \subseteq \Domain{M}$-কে উপগঠনের সংজ্ঞাক্ষেত্র বলা হয়েছে, ফলে $N=\emptyset$-ও অন্তর্ভুক্ত হয়; কিন্তু এই সংস্করণে গঠনের সংজ্ঞাক্ষেত্র সংজ্ঞা অনুসারে অশূন্য। বাংলা পাঠে প্রয়োজনীয় “অশূন্য” শর্তটি যোগ করা হয়েছে।} % BN-SRC-112 TARGET-CORRECTION-END
\end{rem}

% prove something about this? Examples?

\end{document}
